% !TEX program = lualatex
\documentclass[12pt, a4paper]{article}


\usepackage{fontspec}
\setmainfont{Bookman Old Style}%{CMU Serif}

\usepackage[english, russian, ukrainian]{babel}

\usepackage{amsmath, amssymb}
\usepackage{geometry}
\geometry{margin=2.5cm}
\usepackage{hyperref}
\hypersetup{colorlinks=true, linkcolor=blue, urlcolor=blue}
\usepackage{microtype}
\usepackage{parskip}
\usepackage{booktabs}
\usepackage{enumitem}
\usepackage{xcolor}
\usepackage{mdframed}
\usepackage{tabularray}
\UseTblrLibrary{booktabs}
\definecolor{boxbg}{RGB}{245,245,250}
\definecolor{accent}{RGB}{30,80,160}
\definecolor{planckblue}{RGB}{31,119,180}
\definecolor{theoryred}{RGB}{214,39,40}

\usepackage{pgfplots}
\pgfplotsset{compat=1.18}
\usepackage{pgfplotstable}

\usepackage[backend=biber,
            style=phys,         % або numeric / authoryear залежно від вимог
            sorting=none,       % сортування в порядку цитування
            maxnames=4,         % якщо авторів більше 4, скорочувати до et al.
            eprint=true,
            doi=true]{biblatex}

\addbibresource{references.bib}

\newmdenv[backgroundcolor=boxbg, linecolor=accent, linewidth=1pt,
          innerleftmargin=10pt, innerrightmargin=10pt,
          innertopmargin=8pt, innerbottommargin=8pt]{keybox}


\title{\bfseries Конспект: від інфлатона до CMB}
\author{}
\date{}

% ---------------------------------------------------------------
% Planck 2018 binned TT data
% Columns: ell  Dl[muK^2]  err_minus  err_plus
% Source: Planck Collaboration 2018, arXiv:1807.06209
% ---------------------------------------------------------------
\pgfplotstableread{plankdata.txt}\planckdata

% ---------------------------------------------------------------
% LCDM theory curve (CAMB, Planck 2018 best-fit parameters)
% H0=67.36, Ob h^2=0.02237, Oc h^2=0.1200,
% tau=0.0544, As=2.1e-9, ns=0.9649
% ---------------------------------------------------------------
\pgfplotstableread{LCDM.txt}\theorydata

\begin{document}
\maketitle
\tableofcontents
\newpage

%% --------------------------------------------------------
\section{Від рівнянь Айнштайна до рівнянь Фрідмана}
%% --------------------------------------------------------

Відправна точка: FLRW-метрика в \emph{космологічному часі}~$t$
\begin{equation}\label{eq:FLRW}
  ds^2 = -dt^2 + a(t)^2\,\delta_{ij}\,dx^i dx^j
  \qquad (k=0,\text{ плоский Всесвіт})
\end{equation}
є точним розв'язком рівнянь Айнштайна
$G_{\mu\nu} + \Lambda g_{\mu\nu} = 8\pi G\,T_{\mu\nu}$
для ізотропного однорідного ідеального плину
$T^{\mu\nu} = \mathrm{diag}(\rho,\,p,\,p,\,p)$.

Ненульові компоненти тензора Айнштайна:
\begin{equation*}
  G^0{}_0 = -3H^2,
  \qquad
  G^i{}_j = -(2\dot H + 3H^2)\,\delta^i{}_j,
  \qquad
  H \equiv \frac{\dot a}{a}.
\end{equation*}
Підставляючи в рівняння поля, отримуємо \textbf{рівняння Фрідмана}:

\begin{keybox}
\textbf{Перше рівняння Фрідмана} ($(00)$-компонента):
\begin{equation}\label{eq:F1}
  H^2 = \left(\frac{\dot a}{a}\right)^2 = \frac{8\pi G}{3}\,\rho + \frac{\Lambda}{3}
\end{equation}
\textbf{Друге рівняння Фрідмана} ($(ii)$-компонента, «прискорення»):
\begin{equation}\label{eq:F2}
  \frac{\ddot a}{a} = -\frac{4\pi G}{3}\,(\rho + 3p) + \frac{\Lambda}{3}
\end{equation}
\textbf{Рівняння збереження} ($\nabla_\mu T^{\mu\nu} = 0$):
\begin{equation}\label{eq:conservation_low}
  \dot\rho + 3H(\rho + p) = 0
\end{equation}
\end{keybox}

\noindent
\textit{Зауваження:}
\begin{itemize}
  \item Рівняння збереження не є незалежним: воно автоматично слідує
        з двох рівнянь Фрідмана через тотожність Б'янкі
        $\nabla_\mu G^{\mu\nu} = 0$.
  \item Космологічна стала $\Lambda$ еквівалентна ідеальній рідині
        з $\rho_\Lambda = \Lambda/(8\pi G)$, $p_\Lambda = -\rho_\Lambda$,
        тобто $w = -1$.
  \item Відхилення від плоскості ($k = \pm 1$) дає доданок
        $-k/a^2$ у першому рівнянні; Planck обмежує $|\Omega_k| < 0.005$,
        тому $k=0$~--- відмінне наближення.
\end{itemize}

Фонова метрика FLRW у \emph{конформному часі} $\eta$ ($d\eta = dt/a$):
\begin{equation*}
  ds^2 = a(\eta)^2\bigl(-d\eta^2 + dx^2 + dy^2 + dz^2\bigr)
\end{equation*}
де $a(\eta)$~--- масштабний фактор, динаміку якого задають рівняння Фрідмана у конформних змінних (де штрих означає похідну за конформним часом $\partial/\partial\eta$, а $\mathcal{H} \equiv a'/a = aH$ --- конформний параметр Хаббла):
\begin{equation*}
  \mathcal{H}^2 = \frac{8\pi G}{3}\,a^2\bar{\rho} + \frac{\Lambda}{3}\,a^2,
  \qquad
  \mathcal{H}' = -\frac{4\pi G}{3}\,a^2(\bar{\rho} + 3\bar{p}) + \frac{\Lambda}{3}\,a^2.
\end{equation*}
Тут величини $\bar{\rho}(\eta)$ та $\bar{p}(\eta)$ задають ідеально однорідну й ізотропну фонову густину та тиск космічного субстрату відповідно. Комбінація цих рівнянь дає важливе чисто кінематичне співвідношення: $\mathcal{H}' = \mathcal{H}^2 [1 - \frac{3}{2}(1+w)]$.

Залежно від того, яка саме компонента енергії-імпульсу домінує на певному етапі еволюції Всесвіту, змінюється рівняння стану $w = \bar{p}/\bar{\rho}$, що критично визначає кінематику конформного масштабного фактора та поведінку сопутнього горизонту Хаббла $\mathcal{H}^{-1}$:

\begin{center}
\begin{tblr}{
colspec={llll},
row{1} = {c,m, font=\bfseries}
}
\toprule
Епоха & Джерело & $w = \bar{p}/\bar{\rho}$ & $a(\eta)$ \\
\midrule
Інфляція     & Інфлатон                & $\approx -1$ & $\propto -\frac{1}{H\eta} \quad (\eta < 0)$ \\
Радіаційна   & Фотони, нейтрино        & $1/3$        & $\propto \eta$ \\
Матеріальна  & Темна матерія + баріони & $0$          & $\propto \eta^2$ \\
Сьогодні     & Темна енергія           & $\approx -1$ & $\propto -\frac{1}{H\eta} \quad (\text{асимптотика})$ \\
\bottomrule
\end{tblr}
\end{center}

\textit{Методичне зауваження щодо знаку $\eta$ на інфляції:} Оскільки на стадії інфляції Всесвіт розширюється експоненціально, інтеграл конформного часу збігається на верхній границі. Для зручності математичного опису (зокрема, вибору вакууму збурень) початок відліку зміщують так, що інфляція завершується в момент $\eta = 0$. Тому на інфляційній епосі конформний час тече від $-\infty$ до $0$, а сопутній горизонт Хаббла $\mathcal{H}^{-1} = -\eta$ не росте, а стрімко стискається, забезпечуючи причинний зв'язок для масштабів, які сьогодні є мегагалактичними.

%% --- Логічний місток до теорії збурень ---
Рівняння Фрідмана та наведені закони розширення описують ідеалізований, повністю симетричний Всесвіт. Проте наявність реальної великомасштабної структури (галактик та їхніх скупчень) свідчить про те, що реальний космос є неоднорідним. Згідно з концепцією гравітаційної (джинсівської) нестабільності, усі сучасні астрофізичні об'єкти виросли через колапс крихітних початкових флуктуацій.

Математично цей перехід моделюється в рамках лінійної теорії збурень, де повна густина енергії $\rho(\eta, \mathbf{x})$ розкладається на просторово однорідний фон $\bar{\rho}(\eta)$ та малі локальні відхилення $\delta\rho(\eta, \mathbf{x})$:
\begin{equation*}
  \rho(\eta, \mathbf{x}) = \bar{\rho}(\eta) + \delta\rho(\eta, \mathbf{x}),
  \qquad \text{де } |\delta\rho| \ll \bar{\rho}.
\end{equation*}

Контраст густини $\delta \equiv \delta\rho / \bar{\rho}$ динамічно пов'язаний із геометрією простору. Темп розширення фону $a(\eta)$ діє як ефективне «тертя Хаббла», що протидіє силі тяжіння та консервує ріст збурень на радіаційній епосі, дозволяючи їм лінійно зростати ($\delta \propto a$) лише після настання матеріальної ери. Фізичні механізми, які згенерували ці первинні неоднорідності $\delta\rho$ на найраніших етапах, безпосередньо пов'язані з динамікою скалярних полів, що розглядається у наступному розділі.

%% --------------------------------------------------------
\section{Інфлатон і потенціал}
%% --------------------------------------------------------

Найбільш природним і математично узгодженим механізмом генерації початкових неоднорідностей $\delta\rho$ є динаміка гіпотетичного скалярного поля~--- \emph{інфлатона} $\phi(\eta, \mathbf{x})$, яке домінувало в енергетичному балансі раннього Всесвіту. У кожній точці простору-часу стан цього поля характеризується його амплітудою, а динамічна еволюція визначається формою його потенціальної енергії $V(\phi)$.

Фізично цей процес можна уявити як класичне «кочення» поля з локального максимуму потенціалу вниз, аналогічно до руху масивного тіла в полі тяжіння за наявності сильного в'язкого тертя. Завдяки просторовій однорідності фону, на початковому етапі ми можемо знехтувати просторовими градієнтами полях і розглядати лише його фонове значення $\phi(\eta)$.

\begin{keybox}
\textbf{Рівняння руху інфлатона (у космологічному часі $t$):}
\begin{equation}\label{eq:klein_gordon}
  \ddot{\phi} + 3H\dot{\phi} + V'(\phi) = 0
\end{equation}
де $V'(\phi) \equiv dV/d\phi$, а член $3H\dot{\phi}$ описує ефект «хабблівського тертя», зумовлений розширенням самого простору.

\textbf{Тензор енергії-імпульсу скалярного поля} дає такі вирази для фонової густини енергії та тиску:
\begin{equation}\label{eq:inflaton_rho_p}
  \bar{\rho}_\phi = \frac{\dot{\phi}^2}{2} + V(\phi), \qquad \bar{p}_\phi = \frac{\dot{\phi}^2}{2} - V(\phi).
\end{equation}
\end{keybox}

З рівнянь (\ref{eq:inflaton_rho_p}) чітко видно природу інфляційного розширення. Якщо кінетична енергія поля є нехтовно малою порівняно з його потенціальною енергією ($\dot{\phi}^2 \ll V(\phi)$), то ефективне рівняння стану задовольняє умову:
\begin{equation*}
  w_\phi = \frac{\bar{p}_\phi}{\bar{\rho}_\phi} \approx -1.
\end{equation*}
У цьому режимі густина енергії залишається практично сталою ($\bar{\rho}_\phi \approx V(\phi) \approx \text{const}$), що через перше рівняння Фрідмана (\ref{eq:F1}) призводить до сталого параметра Хаббла $H \approx \text{const}$ та експоненційного розширення масштабного фактора:
\begin{equation*}
  a(t) = a_0 \, e^{Ht}.
\end{equation*}

%% --------------------------------------------------------
\subsection{Параметри повільного кочення (slow-roll)}
%% --------------------------------------------------------

Режим, за якого Всесвіт розширюється прискорено, а поле котиться достатньо повільно, щоб забезпечити тривалу інфляцію (щонайменше $N \sim 50$--$60$ $e$-folds), формалізується за допомогою безрозмірних параметрів slow-roll. У загальному випадку їх визначають через динаміку параметра Хаббла (так звані параметри Хаббла--Пригожина):

\begin{keybox}
\textbf{Перший і другий параметри повільного кочення:}
\begin{equation}\label{eq:slow_roll_def}
  \varepsilon \equiv -\frac{\dot H}{H^2} = \frac{\dot\phi^2}{2H^2 M_\mathrm{Pl}^2},
  \qquad
  \eta_\mathrm{H} \equiv -\frac{\ddot{\phi}}{H\dot{\phi}}.
\end{equation}
де $M_\mathrm{Pl} = (8\pi G)^{-1/2}$~--- редукована планківська маса.
\end{keybox}

У наближенні повільного кочення, коли прискоренням поля у рівнянні \eqref{eq:klein_gordon} можна знехтувати ($\ddot{\phi} \ll 3H\dot{\phi}$), воно зводиться до першого порядку: $3H\dot{\phi} \approx -V'(\phi)$. Це дозволяє виразити параметри slow-roll безпосередньо через геометричні характеристики самого потенціалу $V(\phi)$. На відміну від першого розділу, де штрих означав диференціювання за конформним часом, тут і надалі у контексті потенціалу штрихами позначено похідні за скалярним полем $\phi$, тобто $V'(\phi) \equiv dV/d\phi$ та $V''(\phi) \equiv d^2V/d\phi^2$:
\begin{equation}\label{eq:slow_roll_pot}
  \varepsilon \approx \epsilon_V \equiv \frac{M_\mathrm{Pl}^2}{2}\left(\frac{V'}{V}\right)^2,
  \qquad
  \eta \approx \eta_V \equiv M_\mathrm{Pl}^2\,\frac{V''}{V}.
\end{equation}
Умова існування інфляції потребує $\varepsilon < 1$ (що еквівалентно $\ddot{a} > 0$). Для того, щоб цей режим тривав достатньо довго, необхідно виконання нерівностей $\varepsilon \ll 1$ та $|\eta| \ll 1$. Інфляція автоматично завершується в точці потенціалу, де поле прискорюється і перший параметр досягає критичного значення $\varepsilon(\phi_\mathrm{end}) = 1$.



Параметри повільного кочення пов'язують фундаментальну фізику надвисоких енергій із космологічними спостереженнями. У першому порядку теорії збурень спектральні характеристики первинних флуктуацій лінійно виражаються через значення $\varepsilon$ та $\eta$ у момент виходу відповідного масштабу за горизонт:
\begin{equation}\label{eq:ns_r_slow_roll}
  n_s - 1 = -2\varepsilon - \eta, \qquad r = 16\varepsilon,
\end{equation}
де $n_s$~--- спектральний індекс скалярних збурень, а $r$~--- тензорне відношення. Сучасні обмеження на анізотропію CMB (Planck~2018 у комбінації з BICEP/Keck~2021) встановлюють жорсткі межі: $n_s = 0.9649 \pm 0.0042$ та $r < 0.036$. Це дозволяє з високою точністю відсікати класи моделей інфляції, які передбачають занадто крутий нахил потенціалу.

%% --------------------------------------------------------
\subsection{Геометрія потенціалу та обмеження спостережень}
%% --------------------------------------------------------

Важливо розуміти, що спостереження реліктового випромінювання дають нам інформацію не про весь потенціал $V(\phi)$, а лише про його дуже вузький проміжок. Масштаби, які ми фіксуємо сьогодні на картах анізотропії температур CMB, виходили за горизонт Хаббла протягом всього лише $\sim\!7$--$8$ $e$-folds (на загальному тлі у $50$--$60$ $e$-folds від початку процесу).



На сьогодні існує понад 100 теоретичних моделей, які конкурують між собою (модель Старобінського $R^2$, хіггсова інфляція, моделі спонтанного порушення симетрії тощо). Багато з них дають ідеальний фіт для великих масштабів (дані Planck).

Проте, якщо на ділянці потенціалу, що відповідає значно меншим масштабам ($k \gg 1 \, \mathrm{Mpc}^{-1}$), які виходили за горизонт значно пізніше, геометрія потенціалу різко змінюється (наприклад, виникає локальне плато або сходинка), режим повільного кочення тимчасово порушується. Як буде показано далі, такий локальний злом симетрії потенціалу дозволяє суттєво змінити первинний спектр потужності, не руйнуючи при цьому космологічну картину на великих кутах.

%% --------------------------------------------------------
\section{Квантові флуктуації і збурення}
%% --------------------------------------------------------

Найбільш фундаментальним досягненням інфляційної парадигми є те, що вона природним чином пояснює походження первинних неоднорідностей Всесвіту. Поверх класичного однорідного фону скалярного поля $\bar{\phi}(\eta)$ та фонової метрики FLRW завжди існують неминучі квантові флуктуації. У лінійному наближенні теорії збурень їх прийнято розділяти на два абсолютно незалежні класи:
\begin{enumerate}
  \item \textbf{Флуктуації інфлатона} $\delta\phi(\eta, \mathbf{x})$~--- генерують первинні скалярні збурення густини матерії (адіабатичні збурення кривини).
  \item \textbf{Флуктуації метрики} $h_{ij}(\eta, \mathbf{x})$~--- просторово-тензорні збурення, які представляють собою первинні стохастичні гравітаційні хвилі.
\end{enumerate}

%% --------------------------------------------------------
\subsection{Походження скалярних збурень у просторі де-Сіттера}
%% --------------------------------------------------------

Розглянемо квантування легкого безмасового скалярного поля у чисто де-Сіттерівському просторі, де параметр Хаббла є строго сталим ($H = \text{const}$). Переходячи у Фур'є-простір за просторовими координатами $\mathbf{x} \to \mathbf{k}$, рівняння Клейна--Гордона для $k$-ї моди флуктуації $\delta\phi_k(\eta)$ набуває вигляду рівняння руху коливальної системи із загасанням у розширюваному фоні \cite{mukhanov_2005}:
\begin{equation}\label{eq:fourier_field}
  \ddot{\delta\phi}_k + 3H\,\dot{\delta\phi}_k + \frac{k^2}{a^2}\,\delta\phi_k = 0,
\end{equation}
де крапка означає похідну за космологічним часом $t$, а $k \equiv |\mathbf{k}|$~--- комодуляційне хвильове число. Еволюція кожної окремої моди критично dépendent від співвідношення між її фізичною довжиною хвилі $\lambda_{\text{phys}} = a/k$ та космологічним горизонтом Хаббла $R_\mathrm{H} = H^{-1}$ \cite{mukhanov_chibisov_1981}:

\begin{itemize}
  \item \textbf{Субхаббловий режим ($k/a \gg H$):} На надмалих просторових масштабах (всередині горизонту) ефектами розширення простору можна знехтувати. Рівняння (\ref{eq:fourier_field}) зводиться до рівняння стандартного гармонічного осцилятора у пласкому просторі Мінковського. Відповідний вибір початкового квантового стану~--- вакуум Банча--Девіса \cite{bunch_davies_1978}~--- задає нормування амплітуди коливань нульових флуктуацій: $|\delta\phi_k| \sim k^{-1/2} / (a\sqrt{2})$.
  \item \textbf{Суперхаббловий режим ($k/a \ll H$):} Унаслідок стрімкого інфляційного розширення фізичний масштаб флуктуації переростає радіус горизонту Хаббла. Мода «виходить за горизонт», де градієнтний член $k^2/a^2$ стає нехтовно малим, і коливальний режим повністю припиняється \cite{hawking_1982}.
\end{itemize}

У суперхаббловому ліміті ($k \to 0$) загальний розв'язок рівняння (\ref{eq:fourier_field}) містить сталу (заморожену) та затухаючу моди. Амплітуда флуктуації асимптотично виходить на стаціонарне «заморожене» значення \cite{guth_pi_1982}:
\begin{equation*}
  |\delta\phi_k| \to \frac{H}{\sqrt{2k^3}},
\end{equation*}
що дозволяє строго визначити безрозмірний безваріантний за масштабом первинний спектр потужності флуктуацій скалярного поля \cite{mukhanov_chibisov_1981}:
\begin{equation}\label{eq:field_power_spec}
  \mathcal{P}_{\delta\phi}(k) \equiv \frac{k^3}{2\pi^2}|\delta\phi_k|^2 = \left(\frac{H}{2\pi}\right)^2.
\end{equation}

З термодинамічного погляду, вираз (\ref{eq:field_power_spec}) відображає наявність характерної \emph{температури простору де-Сіттера} $T_\mathrm{dS} = H/(2\pi)$ \cite{gibbons_hawking_1977}. Геометричний космологічний горизонт подій обмежує доступну спостерігачеві інформацію, що призводить до квантування вакуумних станів та теплового випромінювання макроскопічних флуктуацій, концептуально аналогічно до випромінювання Гокінга для чорних дір \cite{hawking_1975}.

%% --------------------------------------------------------
\subsection{Перехід до первинних збурень кривини. Зламаний спектр}
%% --------------------------------------------------------

Самі по собі флуктуації поля $\delta\phi$ не є калібрувально-інваріантними. Фізичний зміст мають зумовлені ними первинні \emph{збурення просторової кривини} $\mathcal{R}_k$ на суперхабблових масштабах \cite{bardeen_1980}. Вони лінійно пов'язані з $\delta\phi_k$ через швидкість кочення класичного фону $\dot{\bar{\phi}}$:
\begin{equation}\label{eq:curvature_perturbation}
  \mathcal{R}_k = - \frac{H}{\dot{\bar{\phi}}}\,\delta\phi_k \quad \Rightarrow \quad \mathcal{P}_{\mathcal{R}}(k) = \frac{H^2}{2\pi^2 \dot{\bar{\phi}}^2} \left(\frac{H}{2\pi}\right)^2 = \frac{H^4}{4\pi^2 \dot{\bar{\phi}}^2}.
\end{equation}

Формула (\ref{eq:curvature_perturbation}) є ключовою для аналізу моделей із модифікованою інфляційною кінематикою. У стандартному режимі повільного кочення параметри фону змінюються плавно, забезпечуючи майже плоско-інваріантний спектр Харісона--Зельдовича. Проте, якщо на шляху інфлатона зустрічається локальна аномалія геометрії потенціалу (наприклад, різка сходинка або ділянка ультра-повільного кочення USR), швидкість поля стрімко падає ($\dot{\bar{\phi}} \to 0$).

Оскільки амплітуда збурень обернено пропорційна швидкості кочення ($\mathcal{R}_k \propto 1/\dot{\bar{\phi}}$), флуктуації, які виходять за горизонт саме в цей момент, зазнають різкого локального підсилення. Це приводить до формування \textbf{зламаного спектра потужності} з надлишком амплітуди на фіксованому масштабі $k_\mathrm{break}$, що є центральною ідеєю нашої роботи для узгодження з ранніми галактиками від JWST.

%% --------------------------------------------------------
\subsection{Динаміка «Заморожування» та «Розморожування» мод}
%% --------------------------------------------------------

Процес еволюції первинних збурень підпорядковується чіткій часовій схемі, зумовленій кінематикою горизонту подій:
\begin{enumerate}
  \item \textbf{Заморожування (Епоха інфляції):} Фізичний простір розширюється прискорено ($a \propto e^{Ht}$), тоді як радіус Хаббла $H^{-1}$ залишається практично сталим. Як наслідок, комодуляційні масштаби один за одним покидають горизонт ($\lambda_{\text{phys}} > H^{-1}$). У суперхаббловій зоні збурення кривини $\mathcal{R}_k$ стають константними (заморожуються) \cite{bardeen_1980}, повністю зберігаючи інформацію про локальну мікрофізику потенціалу $V(\phi)$.
  \item \textbf{Розморожування (Постінфляційні епохи):} Після завершення інфляції Всесвіт розширюється сповільнено. За такої кінематики лінійний розмір горизонту Хаббла росте швидше за масштабний фактор. Збурення починають послідовно \textbf{повертатися в горизонт} (розморожуватися) \cite{mukhanov_2005}. При цьому першими всередину горизонту заходять малі масштаби (великі $k$), які вийшли з нього наприкінці інфляції, а значно пізніше — великі масштаби (малі $k$).
\end{enumerate}
Повертаючись в горизонт, кожна мода стає джерелом для класичної еволюції неоднорідностей метрики та гідродинамічних коливань космічної плазми.


%% --------------------------------------------------------
\section{Метрика Фрідмана і збурення}
%% --------------------------------------------------------

Для опису реального неоднорідного Всесвіту, структура якого сформувалася з розглянутих у попередньому розділі квантових флуктуацій $\delta\phi$, необхідно вийти за межі ідеально симетричної фонової геометрії Фрідмана--Леметра--Робертсона--Вокера (FLRW). Малі відхилення метрики простору-часу $\delta g_{\mu\nu}(\eta, \mathbf{x})$ від її фонового значення $\bar{g}_{\mu\nu}(\eta)$ класифікуються відповідно до їхньої поведінки під дією просторових обертань координатної системи.

Згідно з фундаментальною теоремою Стюарта--Вокера \cite{stewart_walker_1974}, будь-яке тензорне збурення на однорідному й ізотропному фоні однозначно розкладається на три \emph{незалежні} компоненти (так званий SVT-розклад):
\begin{equation}\label{eq:svt_decomposition}
  \delta g_{\mu\nu} = \delta g_{\mu\nu}^{(\mathrm{S})} + \delta g_{\mu\nu}^{(\mathrm{V})} + \delta g_{\mu\nu}^{(\mathrm{T})}.
\end{equation}

Кожна з цих компонент еволюціонує у лінійному порядку теорії збурень повністю ізольовано, не взаємодіючи з іншими:
\begin{itemize}
  \item \textbf{Скалярні збурення ($\delta g_{\mu\nu}^{(\mathrm{S})}$):} Параметризуються двома незалежними скалярними функціями простору й часу. Вони безпосередньо динамічно пов'язані з локальними флуктуаціями густини матерії $\delta\rho$. Саме скалярні збурення є причиною гравітаційної нестабільності, яка призводить до формування великомасштабної структури Всесвіту.
  \item \textbf{Векторні збурення ($\delta g_{\mu\nu}^{(\mathrm{V})}$):} Описують вихрові рухи космічного субстрату та моменти імпульсу. Оскільки у стандартній космологічній моделі відсутні джерела, здатні постійно підтримувати векторні поля, ці збурення швидко згасають у процесі розширення Всесвіту пропорційно $a^{-2}$ і не відіграють суттєвої ролі у формуванні структури \cite{mukhanov_2005}.
  \item \textbf{Тензорні збурення ($\delta g_{\mu\nu}^{(\mathrm{T})}$):} Представляють собою вільні поперечні безслідові деформації просторової метрики, які інтерпретуються як первинні стохастичні гравітаційні хвилі. Вони народжуються безпосередньо з квантових флуктуацій метрики під час інфляції і поширюються крізь Всесвіт, практично не взаємодіючи з речовиною.
\end{itemize}

%% --------------------------------------------------------
\subsection{Конформно-ньютонівське калібрування}
%% --------------------------------------------------------

Через загальну коваріантність рівнянь ЗТВ, математичний опис збурень містить фіктивні (калібрувальні) ступені вільності, зумовлені вибором системи координат. Для усунення цих артефактів та фіксації однозначного фізичного змісту ми використовуємо \emph{конформно-ньютонівське калібрування} (також відоме як поздовжнє). У такій системі координат збурена метрика Всесвіту, що враховує скалярні потенціали та тензорні гравітаційні хвилі, записується у вигляді:
\begin{equation}\label{eq:perturbed_metric}
  ds^2 = a(\eta)^2 \Bigl[ -(1 + 2\Phi)\,d\eta^2 + \bigl((1 - 2\Psi)\,\delta_{ij} + h_{ij}\bigr)dx^i dx^j \Bigr],
\end{equation}
де функції $\Phi(\eta, \mathbf{x})$ та $\Psi(\eta, \mathbf{x})$ є безрозмірними калібрувально-інваріантними скалярними потенціалами Бардіна \cite{bardeen_1980}. Кожен із доданків у виразі (\ref{eq:perturbed_metric}) має чітку фізичну інтерпретацію:
\begin{enumerate}
  \item $\Phi$~--- ньютонівський гравітаційний потенціал. Він визначає класичне прискорення пробних частинок і релятивістське зміщення частоти фотонів CMB при їхньому проходженні крізь просторові неоднорідності.
  \item $\Psi$~--- потенціал просторової кривини, що характеризує локальне ізотропне стиснення або розтягнення просторових гіперповерхонь. У межах загальної теорії відносності, за відсутності макроскопічного анізотропного напруження у космічному середовищі, просторовий та часовий потенціали є тотожними: $\Phi = \Psi$.
  \item $h_{ij}$~--- тензорне збурення, яке задовольняє умови поперечності ($\partial^i h_{ij} = 0$) та безслідовості ($h^i{}_i = 0$). Воно описує два незалежні стани поляризації гравітаційних хвиль ($h_+$ та $h_\times$).
\end{enumerate}

%% --------------------------------------------------------
\subsection{Динаміка первинних гравітаційних хвиль}
%% --------------------------------------------------------

Підстановка тензорної частини метрики (\ref{eq:perturbed_metric}) у лінеаризовані рівняння Айнштайна за відсутності джерел анізотропного тиску приводить до бездисипативного хвильового рівняння для кожної Фур'є-моди $h_{ij}(\eta, \mathbf{k})$:
\begin{equation}\label{eq:tensor_wave_eq}
  h_{ij}'' + 2\mathcal{H}h_{ij}' + k^2 h_{ij} = 0,
\end{equation}
де штрих означає диференціювання за конформним часом $\eta$. Структура рівняння (\ref{eq:tensor_wave_eq}) ідентична до рівняння руху для флуктуацій скалярного поля, де роль коефіцієнта згасання (хабблівського тертя) відіграє конформний параметр Хаббла $\mathcal{H} = a'/a$.

Аналогічно до скалярних флуктуацій, квантування нульових коливань метрики під час інфляції у першому порядку теорії збурень визначає первинний спектр потужності тензорних збурень \cite{starobinsky_1979}:
\begin{equation}\label{eq:tensor_power_spec}
  \mathcal{P}_h(k) = \frac{2}{\pi^2 M_\mathrm{Pl}^2} \left(\frac{H}{2\pi}\right)^2 = \frac{H^2}{\pi^2 M_\mathrm{Pl}^2}.
\end{equation}
Амплітуда тензорного спектра $\mathcal{P}_h(k)$ безпосередньо фіксує енергетичний масштаб інфляції ($V \propto H^2$), що робить експериментальний пошук первинних гравітаційних хвиль критично важливим для фізики надвисоких енергій.

Для кількісного зіставлення масштабів тензорних та скалярних збурень використовують \emph{тензорно-скалярне відношення} $r$:
\begin{equation}\label{eq:tensor_scalar_ratio}
  r \equiv \frac{\mathcal{P}_h(k_*)}{\mathcal{P}_{\mathcal{R}}(k_*)} = 16\varepsilon,
\end{equation}
де $\varepsilon$~--- перший параметр повільного кочення, визначений у формулі (\ref{eq:slow_roll_def}), а $k_*$~--- опорний комодуляційний масштаб. Сучасні експерименти з вимірювання поляризації CMB встановлюють жорстку верхню межу $r < 0.036$ \cite{bicep_keck_2021}, що накладає фундаментальні обмеження на амплітуду інфляційного потенціалу.


%% --------------------------------------------------------
\section{Реогрів і народження матерії}
%% --------------------------------------------------------

Після того як інфлатон скочується до мінімуму свого потенціалу $V(\phi)$, період експоненційного розширення простору завершується ($\varepsilon \to 1$). Проте Всесвіт у цей момент перебуває у стані екстремально низької ентропії: він є абсолютно пустим, а вся його енергія зосереджена у формі однорідних класичних коливань поля $\phi$ навколо мінімуму потенціалу. Перехід від цієї «холодної» фази до гарячого Всесвіту, заповненого ультрарелятивістською плазмою у стані теплової рівноваги (що знаменує початок радіаційної епохи), описується теорією \emph{реогріву} (reheating) \cite{linde_1982, albrecht_1982}.

Динаміку цього процесу умовно поділяють на дві принципово різні стадії: допертурбативну (пререогрів) та пертурбативну (лінійний розпад).

%% --------------------------------------------------------
\subsection{Пререогрів та параметричний резонанс}
%% --------------------------------------------------------

На початковому етапі коливання інфлатона $\phi(t)$ навколо мінімуму потенціалу (наприклад, $V(\phi) = \frac{1}{2}m^2\phi^2$) відбуваються з високою амплітудою. Якщо інфлатон квантово-механічно пов'язаний з іншими скалярними полями речовини $\chi$ (наприклад, через член взаємодії $g^2\phi^2\chi^2$), то рівняння руху для Фур'є-мод поля народжуваних часток $\chi_k$ набуває вигляду рівняння Матьє \cite{kofman_1994, kofman_1997}:
\begin{equation}\label{eq:mathieu}
  \ddot{\chi}_k + 3H\dot{\chi}_k + \left[ k_{\mathrm{phys}}^2 + m_\chi^2 + g^2\Phi_0^2\sin^2(mt) \right]\chi_k = 0,
\end{equation}
де $\Phi_0(t)$~--- амплітуда коливань інфлатона, яка поступово затухає через розширення Всесвіту.

Періодична зміна коефіцієнтів у рівнянні (\ref{eq:mathieu}) призводить до ефекту \emph{параметричного резонансу} у певних частотних діапазонах (зонах нестабільності). Цей процес отримав назву \textbf{пререогріву} (preheating). Замість повільного поштучного народження часток, тут відбувається експоненціально швидке, вибухове некогерентне народження квантів поля $\chi$ (так зване явище стимульованого випромінювання Бозе--Айнштайна). За лічені осциляції класичне поле інфлатона фрагментується, а густина енергії часток $\chi$ стрімко зростає.

%% --------------------------------------------------------
\subsection{Пертурбативний розпад та термалізація}
%% --------------------------------------------------------

Коли амплітуда коливань падає, параметричний резонанс вимикається, і подальша еволюція переходить у стадію повільного пертурбативного розпаду. Математично цей процес моделюється введенням феноменологічного члена загасання~--- швидкості розпаду інфлатона $\Gamma_\phi$~--- безпосередньо у рівняння балансу енергії для густини інфлатона $\rho_\phi$ та народжуваної релятивістської плазми (радіації) $\rho_r$:
\begin{align}
  \dot{\rho}_\phi + 3H\rho_\phi &= -\Gamma_\phi \rho_\phi, \label{eq:reheat_rho_phi} \\
  \dot{\rho}_r + 4H\rho_r &= +\Gamma_\phi \rho_\phi. \label{eq:reheat_rho_r}
\end{align}

Усереднене за часом рівняння стану коливного інфлатона на дні параболічного потенціалу відповідає пилоподібній матерії ($w_\phi = 0$). Тому доки $H \gg \Gamma_\phi$, Всесвіт розширюється за законом матеріальної епохи ($a \propto t^{2/3}$), а густина енергії плазми підживлюється розпадом поля.



Процес закінчується, коли темп розширення Всесвіту зрівнюється зі швидкістю розпаду частинок: $H \sim \Gamma_\phi$. У цей момент залишки інфлатона майже миттєво розпадаються, а релятивістські продукти розпаду за рахунок пружних та непружних процесів розсіяння зазнають швидкої \emph{термалізації}, встановлюючи єдину температуру середовища.

\begin{keybox}
\textbf{Температура реогріву Всесвіту ($T_{\mathrm{rh}}$):}
Визначається з умови $H^2 = \frac{8\pi G}{3}\rho_r$ при $H \approx \Gamma_\phi$:
\begin{equation}\label{eq:t_reheating}
  T_{\mathrm{rh}} = \left( \frac{90}{8\pi^3 g_*} \right)^{1/4} \sqrt{M_{\mathrm{Pl}}\Gamma_\phi},
\end{equation}
де $g_* \sim 100$~--- ефективна кількість релятивістських ступенів вільності Стандартної моделі при надвисоких енергіях.
\end{keybox}

Значення $T_{\mathrm{rh}}$ є фундаментальним параметром для подальшої гарячої історії Всесвіту. Воно задає початкову температуру для баріогенезису, фазових переходів та термодинамічного загартування темної матерії. Що найважливіше для нашої моделі збурень — саме в цей момент гаряча фотон-баріонна плазма «успадковує» просторові неоднорідності кривини $\mathcal{R}_k$, які були згенеровані під час інфлатонного зламу та «заморожені» на суперхабблових масштабах.

%% --------------------------------------------------------
\section{Гідродинаміка фотон-баріонної плазми та акустичні осциляції}
%% --------------------------------------------------------

Після завершення реогріву Всесвіт вступає у радіаційно-доміновану епоху. Через екстремально високі температури баріонна матерія (протони та електрони) повністю іонізована і за рахунок інтенсивного томсонівського розсіяння фотонів на вільних електронах міцно пов'язана з випромінюванням. Ця система моделюється як єдина релятивістська \emph{фотон-баріонна плазма}, яка еволюціонує у полі ньютонівських гравітаційних потенціалів $\Phi$ та $\Psi$, успадкованих від інфляційної епохи \cite{mukhanov_2005}.

Динаміка скалярних збурень у такому середовищі визначається конкуренцією двох сил: гравітаційного притягання, яке намагається стиснути речовину у потенціальні ями, та внутрішнього ізотропного тиску фотонного газу, що чинить опір стисненню.

%% --------------------------------------------------------
\subsection{Рівняння акустичного осцилятора}
%% --------------------------------------------------------

У лінійному порядку теорії збурень гідродинамічні рівняння неперервності та Ейлера для Фур'є-моди фотонного контрасту густини $\delta_\gamma \equiv \delta\rho_\gamma / \bar{\rho}_\gamma$ на масштабах, значно менших за горизонт ($k\eta \gg 1$), зводяться до рівняння вимушених коливань із тертям \cite{dodelson_2020}:
\begin{equation}\label{eq:acoustic_oscillator}
  \frac{d}{d\eta}\left[ (1 + R)\,\delta_\gamma' \right] + \frac{k^2}{3}\,\delta_\gamma = -\frac{k^2}{3}(1 + R)\,\Phi,
\end{equation}
де штрих означає похідну за конформним часом $\eta$, а безрозмірний параметр $R(\eta) \equiv \frac{3\bar{\rho}_b(\eta)}{4\bar{\rho}_\gamma(\eta)} \propto a(\eta)$ описує відносну інерційну вагу баріонної компоненти.

Швидкість поширення звуку у такому середовищі є релятивістською і динамічно змінюється внаслідок збільшення маси баріонів при розширенні простору:
\begin{equation}\label{eq:sound_speed}
  c_s(\eta) = \frac{1}{\sqrt{3(1 + R(\eta))}}.
\end{equation}

Рівняння (\ref{eq:acoustic_oscillator}) є математичним описом \textbf{космічного акустичного осцилятора}. На субхабблових масштабах кожна окрема мода $k$ починає коливатися (стискатися та розширюватися) навколо зміщеного гравітаційним потенціалом $\Phi$ положення рівноваги. Оскільки всі моди з однаковим хвильовим числом $k$ були згенеровані інфляцією когерентно (тобто мали однакову початкову фазу), їхні коливання відбуваються строго синхронно в часі.

%% --------------------------------------------------------
\subsection{Баріонні акустичні осциляції (BAO)}
%% --------------------------------------------------------

Коливання фотон-баріонної плазми тривають до моменту \emph{рекомбінації гідрогену} ($\eta = \eta_{\mathrm{rec}}$), коли температура Всесвіту падає до $\sim 3000$~К. Вільні електрони об'єднуються з протонами, фотони миттєво перестають розсіюватися і починають вільно поширюватися у просторі, утворюючи сучасну картину кутової анізотропії CMB \cite{peebles_yu_1970}.

Максимальна відстань, яку звукова хвиля встигає пройти у плазмі від початку гарячої епохи до моменту рекомбінації, фіксує фундаментальний лінійний масштаб у космології~--- \emph{радіус звукового горизонту} $r_s$:
\begin{equation}\label{eq:sound_horizon}
  r_s(\eta_{\mathrm{rec}}) = \int_0^{\eta_{\mathrm{rec}}} c_s(\eta)\,d\eta \approx 147 \text{ Mpc}.
\end{equation}

Ті комодуляційні моди $k$, які в момент рекомбінації підійшли до точки свого максимального стиснення або розширення, відображаються на кутовому спектрі потужності у вигляді серії акустичних піків.

\begin{figure}[htbp]
\centering
\begin{tikzpicture}
  \begin{axis}[
    width=0.85\textwidth,
    height=5.5cm,
    xmin=0, xmax=250,
    ymin=-0.05, ymax=1.1,
    xlabel={Відстань від первинного збурення, $r$ [Мпк]},
    ylabel={Контраст густини баріонів, $\Delta\rho_b/\bar{\rho}_b$},
    grid=both,
    grid style={line width=.1pt, color=gray!20},
    major grid style={line width=.2pt, color=gray!40},
    tick label style={font=\footnotesize},
    label style={font=\small},
    legend style={at={(0.95,0.95)}, anchor=north east, font=\footnotesize, cells={anchor=west}},
    axis lines=left,
    clip=false
  ]
    % Моделювання профілю BAO (сферичний сплеск + дифузійне затухання Шовка)
    \addplot[
      domain=0:240,
      samples=200,
      color=accent,
      line width=1.5pt
    ] {exp(-0.015*x) + 0.35*exp(-((x-147)/12)^2)};
    \addlegendentry{Профіль густини речовини}

    % Позначення радіуса звукового горизонту
    \draw[dashed, color=theoryred, line width=1pt] (axis cs:147,-0.05) -- (axis cs:147,0.9);
    \node[anchor=south, color=theoryred, font=\footnotesize] at (axis cs:147,0.9) {$r_s \approx 147\text{ Мпк}$};

    % Позначення дифузійного затухання
    \node[anchor=west, color=gray, font=\scriptsize] at (axis cs:165,0.2) {Затухання Шовка};
    \draw[->, color=gray, >=stealth] (axis cs:162,0.18) -- (axis cs:155,0.1);
  \end{axis}
\end{tikzpicture}
\caption{Профіль поширення первинного адіабатичного збурення у баріонній компоненті плазми. Чітко виражений надлишок речовини на радіусі звукового горизонту $r_s$ відповідає сучасному масштабу BAO.}
\label{fig:bao_profile}
\end{figure}

Проте акустичні хвилі залишають свій слід не лише у випромінюванні. Після рекомбінації фотони пішли геть, а баріонна речовина під дією сили тяжіння залишилася «замороженою» у вигляді сферичної оболонки радіусом $r_s$ навколо кожної початкової неоднорідності темної матерії. Цей феномен отримав назву \textbf{баріонних акустичних осциляцій (BAO)} \cite{eisenstein_2005}.

У сучасну епоху масштаб BAO проявляється як чітко фіксований надлишок у просторовій двоточковій кореляційній функції розподілу галактик на відстані $\sim 147$~Мпк. Оскільки цей лінійний розмір жорстко заданий лінійною фізикою раннього Всесвіту (\ref{eq:sound_horizon}), BAO слугує ідеальною «стандартною лінійкою» (standard ruler). Вимірювання цього масштабу на різних червоних зміщеннях за допомогою великомасштабних оглядів неба дозволяє з високою точністю реконструювати темп розширення Всесвіту $H(z)$ та геометричні параметри темної енергії, незалежно від даних реліктового випромінювання.

%% --------------------------------------------------------
\section{Ефект Закса--Вольфа та формування спектра CMB}
%% --------------------------------------------------------

Кутова анізотропія температури реліктового випромінювання (CMB), що спостерігається нами сьогодні, є безпосереднім відбитком («знімком») стану акустичних осциляцій фотон-баріонної плазми в епоху рекомбінації Всесвіту ($z_{\mathrm{rec}} \approx 1100$). Флуктуації температури у заданому напрямку на небі $\hat{\mathbf{n}}$ розкладаються за сферичними гармоніками, а їхні статистичні властивості повністю описуються кутовим спектром потужності $\mathcal{D}_\ell^{TT} \equiv \frac{\ell(\ell+1)}{2\pi}C_\ell^{TT}$, де мультиполь $\ell$ обернено пропорційний кутовому масштабу $\theta \sim 180^\circ / \ell$ \cite{dodelson_2020}.

Спостережувані варіації температури $\delta T / T$ формуються внаслідок дії трьох основних фізичних механізмів на поверхні останнього розсіяння (LSS) \cite{hu_sugiyama_1995}:
\begin{equation}\label{eq:cmb_sources}
  \frac{\delta T}{\text{T}}(\hat{\mathbf{n}}) = \left( \frac{1}{4}\delta_\gamma + \Phi \right)_{\mathrm{LSS}} + \left( \hat{\mathbf{n}} \cdot \mathbf{v}_b \right)_{\mathrm{LSS}} + \int_{\eta_{\mathrm{LSS}}}^{\eta_0} \left( \Phi' + \Psi' \right) d\eta.
\end{equation}

Кожен із доданків у виразі (\ref{eq:cmb_sources}) відповідає за свій масштаб на спектрі:
\begin{enumerate}
  \item \textbf{Звичайний ефект Закса--Вольфа ($\frac{1}{4}\delta_\gamma + \Phi$):} Домінує на великих кутових масштабах (малі мультиполі $\ell \lesssim 100$), які на момент рекомбінації ще перебували поза горизонтом Хаббла. Він поєднує в собі власну термодинамічну анізотропію фотонів та гравітаційне червоне зміщення, якого фотони зазнають при виході з потенціальних ям $\Phi$. Для адіабатичних збурень $\frac{1}{4}\delta_\gamma = -\frac{2}{3}\Phi$, що дає сумарний ефект $\frac{1}{3}\Phi$.
  \item \textbf{Доплерівський ефект ($\hat{\mathbf{n}} \cdot \mathbf{v}_b$):} Зумовлений кінематичним рухом (швидкістю $\mathbf{v}_b$) фотон-баріонної плазми у полі збурень. Цей внесок зміщений за фазою відносно ефекту Закса--Вольфа і частково заповнює мінімуми між акустичними піками.
  \item \textbf{Інтегральний ефект Закса--Вольфа (ISW, інтегральний член):} Виникає, коли фотони CMB по дорозі до спостерігача пролітають крізь гравітаційні потенціали, що змінюються у часі ($\Phi' \neq 0$). Ранній ISW проявляється відразу після рекомбінації через залишковий вплив радіації, а пізній ISW~--- на надвеликих масштабах ($\ell < 10$) в сучасну епоху через прискорене розширення Всесвіту під дією темної енергії.
\end{enumerate}

%% --------------------------------------------------------
\subsection{Акустичні піки та затухання Шовка}
%% --------------------------------------------------------

На субхабблових масштабах ($\ell \gtrsim 100$) когерентні звукові коливання формують чітку серію \emph{акустичних піків}. Перший і найвищий пік при $\ell \approx 220$ ($\theta \sim 1^\circ$) відповідає моді, яка встигла зазнати рівно одного максимального стиснення у потенціальній ямі до моменту рекомбінації. Його положення є чутливим індикатором просторової геометрії Всесвіту \cite{planck_2018}. На мультиполях $\ell \gtrsim 1000$ спектр експоненціально загасає через ефект \emph{дифузійного затухання} (затухання Шовка)~--- фотони мають скінченну довжину вільного пробігу і в процесі дифузії просто замивають дрібномасштабні флуктуації температури \cite{silk_1968}.

\begin{figure}[b!]
\centering
\begin{tikzpicture}
\begin{semilogxaxis}[
    width  = 14.5cm,
    height = 8.5cm,
    xlabel = {Мультипольний момент $\ell$},
    ylabel = {$\mathcal{D}_\ell^{TT}\ [\mu\mathrm{K}^2]$},
    xmin   = 2,   xmax = 2600,
    ymin   = 0,   ymax = 6800,
    xtick  = {2, 10, 30, 100, 300, 1000, 2500},
    xticklabels = {2, 10, 30, 100, 300, 1000, 2500},
    minor xtick = {5,20,50,200,500,2000},
    ytick  = {0, 1000, 2000, 3000, 4000, 5000, 6000},
    grid   = both,
    grid style = {line width=0.3pt, draw=gray!25},
    major grid style = {line width=0.5pt, draw=gray!40},
    tick label style = {font=\small},
    label style      = {font=\small},
    legend style = {
        at={(0.6,0.97)}, anchor=north east,
        font=\small, fill=white, fill opacity=0.9,
        draw=gray!60, inner sep=4pt, row sep=1pt
    },
    clip = true,
]

%--- Theory curve (LCDM, CAMB) ---
\addplot[
    color     = theoryred,
    line width = 1.6pt,
    smooth,
] table[x=ell, y=Dl] {\theorydata};
\addlegendentry{$\Lambda$CDM (CAMB, Planck\,2018 best-fit)}

%--- Planck 2018 data with error bars ---
\addplot[
    only marks,
    mark       = *,
    mark size  = 1.4pt,
    color      = planckblue,
    error bars/.cd,
        y dir  = both,
        y explicit,
] table[
    x    = ell,
    y    = Dl,
    y error minus = errm,
    y error plus  = errp,
] {\planckdata};
\addlegendentry{Planck\,2018 (збіновані дані TT)}

%--- Annotate acoustic peaks ---
\node[font=\tiny, text=gray!70, above] at (axis cs:220,5732) {1-й пік};
\node[font=\tiny, text=gray!70, above] at (axis cs:540,2592) {2-й пік};
\node[font=\tiny, text=gray!70, above] at (axis cs:810,2540) {3-й пік};

%--- Sachs-Wolfe plateau annotation ---
\draw[->, gray!60, thin] (axis cs:6,1300) -- (axis cs:30,1070)
    node[midway, left, font=\tiny, text=gray!60] {плато Закса--Вольфа};

\end{semilogxaxis}
\end{tikzpicture}
\caption{Кутовий TT-спектр потужності CMB. Точки~--- дані Planck~2018
         \cite{planck2018params}.
         Крива~--- базова модель $\Lambda$CDM, розрахована CAMB
         з параметрами таблиці~\ref{tab:params}.
         Видно три акустичні піки та плато Закса--Вольфа на великих
         кутових масштабах ($\ell \lesssim 30$).}
\label{fig:cmb_tt}
\end{figure}


%% --------------------------------------------------------
\section{Поляризація CMB}
%% --------------------------------------------------------

%% --------------------------------------------------------
\subsection*{Механізм поляризації: розсіяння Томсона}
%% --------------------------------------------------------

Фотони поляризуються при останньому розсіянні на електронах.
Томсонівське розсіяння поляризує фотон \emph{лише} якщо навколо електрона
існує \textbf{квадрупольна} температурна анізотропія: якщо температура
однакова з усіх боків~--- поляризації немає.

Безпосередньо до рекомбінації фотонів і електронів ще достатньо для
ізотропізації, тому квадруполь малий і поляризація слабка.
Але під час самої рекомбінації ($\Delta z \sim 80$) оптична глибина
падає швидко~--- квадруполь встигає вирости і поляризація «відбивається»
в поверхні останнього розсіяння.

Напрямок поляризації залежить від \textbf{поперечного} перепаду температури
навколо точки розсіяння.

%% --------------------------------------------------------
\subsection*{E- і B-моди}
%% --------------------------------------------------------

Поляризаційне поле $P_{ab}(\hat n)$ на сфері розкладають на два типи
(аналогічно до розкладу вектора на градієнт і ротор):

\begin{center}
\begin{tblr}{XXX}
\toprule
Тип & Джерело & Статус \\
\midrule
E-моди (градієнтний, симетричний візерунок) & Скалярні флуктуації & Виміряні \\
B-моди (ротаційний, закручений візерунок)  & Тензорні флуктуації (GW) & Поки не знайдені \\
\bottomrule
\end{tblr}
\end{center}

\begin{keybox}
\textbf{Чому скалярні збурення не дають B-мод?}
Скалярні збурення мають потенціальний (паритетно-парний) характер:
$\Phi(\mathbf{x})$ визначає поле швидкостей $\mathbf{v} = \nabla\Phi$.
E-моди мають ту ж симетрію паритету~--- вони парні.
B-моди псевдоскалярні (непарні), і лінійний скалярний збурення їх не генерує.
Тензорне збурення $h_{ij}$ має обидва стани, в тому числі непарний~---
звідси B-моди від гравітаційних хвиль.
\textit{Практичний забруднювач:} лінзування на великих масштабах перетворює
E-моди в B-моди~--- треба враховувати при пошуку первинних B-мод.
\end{keybox}

Знайти B-моди~--- означає отримати прямий доказ первинних гравітаційних хвиль
і квантової природи гравітації в режимі малих збурень.
Поточна верхня межа: $r < 0.036$ (95\% CL, BICEP/Keck~2021).

%% --------------------------------------------------------
\section{Хаотична інфляція Лінде}
%% --------------------------------------------------------

Лінде (1983): початкове значення $\phi$ хаотичне~--- різне в різних точках простору.
Найпростіший потенціал: $V(\phi) = \frac{1}{2}m^2\phi^2$.

Де $\phi$ велике~--- інфляція іде довго, виникає величезна «бульбашка».
Наш Всесвіт~--- одна з таких бульбашок. Інші бульбашки недосяжні~---
простір між ними розширюється швидше за світло.
Це \textbf{вічна інфляція} і мультивсесвіт як математичний наслідок.

%% --------------------------------------------------------
\section{Наша задача: зламаний спектр і JWST}
%% --------------------------------------------------------

%% --------------------------------------------------------
\subsection{Мотивація}
%% --------------------------------------------------------

JWST знайшов масивні галактики при $z \sim 10$--$12$, яких у стандартній
моделі бути не повинно~--- не вистачало часу утворитись.

%% --------------------------------------------------------
\subsection{Ідея}
%% --------------------------------------------------------

Якщо потенціал інфлатона має \textbf{сходинку}, м'яч сповільнюється~---
$\dot{\phi} \to 0$~--- і флуктуації що заморожуються в цей момент мають
більшу амплітуду. Зламаний спектр:
\begin{equation*}
  P(k) = A_s \left(\frac{k}{k_*}\right)^{n_s-1}
  \times \mathcal{F}(k;\, k_\mathrm{break},\, \Delta n)
\end{equation*}
Більше потужності на малих масштабах $k \gtrsim 1\,\mathrm{Mpc}^{-1}$~---
гало тієї самої маси утворюються \textbf{раніше}~--- узгоджується з JWST.

%% --------------------------------------------------------
\subsection{Чому не суперечить CMB і $\sigma_8$}
%% --------------------------------------------------------

Planck чутливий до $k \lesssim 0.2\,\mathrm{Mpc}^{-1}$~--- злом не бачить.
$\sigma_8$ визначається масштабом $\sim 8\,\mathrm{Mpc}/h$, тобто
$k \sim 0.1$--$0.5\,\mathrm{Mpc}^{-1}$~--- майже не змінюється.

%% --------------------------------------------------------
\subsection{План роботи}
%% --------------------------------------------------------

\begin{enumerate}
  \item Огляд літератури: чи зроблено вже саме це поєднання.
  \item Мінімальна модель: 2 нових параметри $(k_\mathrm{break},\, \Delta n)$.
  \item CAMB/CLASS: розрахунок $P(k)$, $C_\ell$, $\sigma_8$.
  \item HMF (\texttt{hmf}, Python): порівняння з каталогами JWST.
  \item MCMC (Cobaya): posterior по повному набору параметрів.
  \item Стаття: чи є сумісна область параметрів для Planck + JWST + $\sigma_8$.
\end{enumerate}


\printbibliography

\end{document}
